\section{Introduction}
\label{sec:intro}

The deployment of specialized, large-scale deep learning models across multiple distinct task domains represents a foundational pillar of modern machine learning services. To mitigate the prohibitive computational and memory footprints associated with hosting individual full-parameter models, the field has increasingly adopted Parameter-Efficient Fine-Tuning (PEFT) frameworks, with Low-Rank Adaptation (LoRA) \cite{hu2021lora} serving as the dominant standard. By freezing a massive shared base backbone and learning small, task-specific parameter matrices, LoRA dramatically reduces storage requirements. However, in dynamic, real-world edge serving environments where heterogeneous request streams from diverse task domains arrive concurrently, a critical system challenge emerges: how to orchestrate these specialized expert adapters dynamically on a sample-by-sample basis without incurring unsustainable inference latencies or memory buffers.

Dynamic ensembling and weight-merging frameworks have arisen to address this challenge by combining task-specific experts on-the-fly. However, current state-of-the-art routing techniques are constrained by a fundamental dilemma known as the \emph{Routing Paradox} \cite{sable2025}. To route an incoming query to the correct task-specific adapter, the system must evaluate a routing function based on the input's features. If this routing function relies on mid-to-late layer representations, the system cannot adapt the early layers of the network without executing the base model twice (once for routing feature extraction, and once for specialized forward propagation), which doubles the serving latency. 

To resolve the Routing Paradox, existing non-parametric ensembling methods such as SABLE \cite{sable2025} and SPS-ZCA \cite{spszca2025} resort to what we term the \textbf{Early-Feature Loss Trade-Off}. These methods enforce "Late Adaptation," freezing and leaving the first several layers (often up to 10 out of 12 layers) completely unadapted. Consequently, any task-specific specialized features learned in the early-to-mid layers of the experts are discarded. Our empirical evaluations demonstrate that this early-layer restriction severely limits the representation capacity of the merged model. Under highly overlapping task representation spaces, SABLE's late adaptation constraints severely restrict its performance, achieving a Joint Mean accuracy of only 55.30\%---falling short of the Expert Ceiling (69.72\%) and highlights the significant capacity constraints of freezing early-to-mid blocks.

Alternatively, systems can employ classical parametric gating networks (e.g., linear classification routers) trained on low-data calibration splits ($B_{\text{cal}} = 64$ per task) \cite{pfsr2024}. While these parametric routers appear highly performant under large, homogeneous batches, they suffer from a failure mode under realistic streaming scenarios: \textbf{Vectorization Collapse}. Because they are optimized under batch-average objectives, their decision boundaries overfit to the calibration distribution. When deployed under heterogeneous, vectorized stream regimes with a batch size of $B=1$ (batch-independent serving), their routing accuracy degrades, plunging joint accuracy to 52.36\%. Other methods attempt to bypass this via complex scheduling frameworks such as Micro-Batch Homogenization (MBH) \cite{pfsr2024}, but these introduce heavy systems-level state buffers and $O(K)$ sequential execution passes, destroying edge serving guarantees.

In this work, we propose **PEAR (Patch-Embedding Activation Routing)**, a parameter-free, closed-form routing framework that relentlessly applies Occam's razor to resolve these structural limitations. PEAR argues that the base model's first structural projection—the frozen Patch Embedding layer (Layer 0)—already maps raw input pixels into an aligned, highly informative representation space at virtually zero extra compute. By computing routing coefficients immediately at Layer 0, PEAR completely bypasses the Routing Paradox, allowing task-specific LoRA adapters to be activated and blended across \textbf{100\% of the network depth} (all 12 layers adapted). 

PEAR achieves high-performance ensembling through a sequence of minimalist, non-parametric calibration steps:
\begin{enumerate}
    \item \textbf{Zero-Shot Patch Centroids (ZPC):} Extracting representative class coordinates in the Layer 0 space using a tiny, weakly-supervised calibration set, requiring zero additional gradients or optimization parameters.
    \item \textbf{Unit-Norm Cosine Projection:} Evaluating cosine similarity on a unit-norm hypersphere to ensure strict scale invariance across diverse visual manifolds.
    \item \textbf{Intra-Task Dispersion Calibration (IDC):} Normalizing similarity scores by expected in-distribution calibration variance to resolve asymmetric representation density.
\end{enumerate}

By feeding the calibrated similarities into a temperature-scaled Softmax, PEAR obtains robust, sample-specific routing weights at Layer 0. These are then used to guide Single-Pass Activation Blending (SPS) dynamically across all layers in a single forward pass, maintaining flat $O(1)$ latency and zero dynamic state memory. To evaluate PEAR under precise representation sharing conditions, we construct a high-fidelity 12-layer synthetic Vision Transformer representation sandbox designed in PyTorch. Tasks are simulated via 1D vector subspaces modeled after standard visual datasets of varying difficulty (MNIST, Fashion-MNIST, CIFAR-10, and SVHN) under a highly challenging overlapping task subspace layout. Crucially, as an intentional methodological stress-test to evaluate ensembling robustness under severe noise and degraded conditions, the SVHN task in our sandbox is configured with a low expert classification ceiling of $19.68\%$, while clean tasks such as MNIST maintain a $100.00\%$ expert ceiling. Across extensive evaluations on this representation sandbox, PEAR achieves a Joint Mean accuracy of \textbf{59.34\%} across all batch deployments. This completely eliminates Vectorization Collapse, outperforming Static Uniform Weight Merging by \textbf{+8.70\%} and the Linear Router by \textbf{+6.98\%--+8.52\%}, while significantly outperforming the late adaptation of SABLE SOTA (\textbf{55.30\%}) by a substantial \textbf{+4.04\%} absolute accuracy gain.

In summary, our key contributions are:
\begin{itemize}
    \item We identify and formalize the \emph{Early-Feature Loss Trade-Off} in late-adaptation non-parametric routers and \emph{Vectorization Collapse} in low-data parametric routing.
    \item We present \textbf{PEAR}, a training-free, non-parametric, closed-form ensembling framework that performs high-fidelity routing inside the Patch Embedding layer (Layer 0).
    \item We demonstrate that PEAR achieves state-of-the-art dynamic ensembling performance across homogeneous and heterogeneous stream configurations (with $B=256$ and $B=1$) inside a high-fidelity 12-layer synthetic representation sandbox designed in PyTorch.
    \item We bridge the simulation-to-real-world gap by evaluating PEAR on actual real-world images from MNIST, Fashion-MNIST, CIFAR-10, and SVHN using a pre-trained $\mathtt{vit\_tiny\_patch16\_224}$ backbone. We empirically demonstrate that shifting the routing boundary slightly deeper (to Layer 1 or 2) resolves the Global-Average-Color Routing Paradox, achieving up to \textbf{95.31\%} routing accuracy, outperforming trained pre-backbone CNN routers (\textbf{91.02\%}) while preserving over 83\% of the network's adaptability with minimal compute overhead.
\end{itemize}
